\section{Problem Formulation}
\label{sec:formulation}

We formalize the workflow as it is implemented; no step below is learned from data.

\paragraph{Submissions.} For a monitored organization (entity) $e$ and assessment
period $t$, a \emph{submission} $S_{e,t}$ is a versioned set of records partitioned
into six evidence categories (alerts, cases, investigation steps, escalations,
dispositions, assets). Each record $r$ carries its source provenance (artifact digest
and record locator). Validation is a total function
$V(S_{e,t})\in\{\mathit{accept},\mathit{reject}\}$ that parses every record, resolves
references (e.g.\ a step must reference an existing case) and enforces chronology; it
is all-or-nothing, so a single invalid record rejects the version.

\paragraph{Findings.} Each analytical worker $w_j$, $j=1,\dots,16$, maps
an accepted submission (and, for peer analysis, a declared cohort) to findings
\begin{equation}
f=(\rho_f,\ \mathit{subject}_f,\ E_f,\ s_f,\ \theta_f,\ c_f),
\end{equation}
with rule $\rho_f$, subject, evidence references $E_f\subseteq S_{e,t}$, statistic $s_f$,
threshold $\theta_f$ and confidence $c_f\in[0,1]$. A finding is a \emph{signal} only if
$E_f\neq\emptyset$; findings without evidence references are withheld. Let $F_e$ be the
signal findings of the entity's run.

\paragraph{Rule families.} Most execution-gap rules are \emph{existence rules}: they
fire when at least one qualifying record exists, e.g.\ a critical alert closed faster
than a severity-specific limit. Peer deviation compares an entity metric $x_e$ with the
median $\tilde m$ and median absolute deviation $\mathrm{MAD}$ of a declared cohort
$C$ (after symmetric 10\% trimming when $|C|\ge 4$) and fires when
\begin{equation}
|C|\ge 3 \quad\text{and}\quad |x_e-\tilde m| \ge 2\,\mathrm{MAD}(C);
\label{eq:peer}
\end{equation}
within-entity anomalies use the same statistic with factor $3$ over the entity's own
values. MAD is unscaled.

\paragraph{Risk fusion.} A fixed map $\delta$ assigns each rule to one of
$7$ dimensions $d$ with weights $w_d$ (execution gap
25, peer deviation 20, detection gap
15, negative space 15, anomaly
15, investigation quality 5,
escalation discipline 5; sum 100). With
$C_d=\sum_{f\in F_e,\ \delta(\rho_f)=d} c_f$,
\begin{equation}
R_d = w_d\,\frac{C_d}{1+C_d},\qquad R_e=\min\Bigl(100,\ \sum_d R_d\Bigr).
\label{eq:risk}
\end{equation}
Each dimension saturates at its weight. Two properties matter later: $R_e$ is a
heuristic score, not a probability or a severity; and it depends on summed confidences,
not on how \emph{prevalent} a problem is within the submission. Findings from different
families naming the same subject are reported as corroborating clusters that do not
change $R_e$.

\paragraph{Prioritization.} Entities are ordered by the priority score
\begin{equation}
P_e = R_e + 10\,\kappa_e + 5\,\tau_e + 2\,h_e,
\label{eq:priority}
\end{equation}
where $\kappa_e\in\{0,\tfrac14,\tfrac35,1\}$ encodes the confidence bucket of the entity's
findings, $\tau_e\in[0,1]$ decays linearly with run age over 30 days, and $h_e$ counts
high-severity signals. The coefficients are fixed defaults, not fitted.

\paragraph{Decision and binding.} A supervisor (or organization administrator)
records a decision $D_e$ on the reviewed state. Finalization builds a canonical
document $\mathit{doc}_e$ (decision, findings, observations, risk, recommendations,
submission records and their provenance), computes $h_e=\mathrm{SHA3\text{-}256}(\mathit{doc}_e)$,
signs $h_e$ with ML-DSA-65, stores a receipt, and appends a ledger entry
$\ell_i$ whose hash is $\mathrm{SHA3\text{-}256}$ of its canonical body including the
predecessor hash $\ell_{i-1}$. Verification rebuilds $\mathit{doc}_e$ from the live
records and checks the digest, the decision binding, the signature and the chain.

\paragraph{Questions.} Given this pipeline, we ask whether the declared mechanisms
behave as specified under controlled conditions, how they compare with simple
baselines, how they respond to imperfect evidence and cohort choice, what orchestration
and integrity cost, and whether $R_e$ tracks an independent outcome on real data
(Section~\ref{sec:method}).
